\section{Conclusão e trabalhos futuros}

\subsection{Síntese das contribuições}

Neste trabalho, foi desenvolvida uma prova de conceito reprodutível de \textit{chatbot} legislativo com RAG voltada ao acervo de discursos da 56\textordfeminine{} Legislatura do Senado Federal. A solução compreende um \textit{pipeline} de preparação documental, uma base vetorial, uma interface de consulta e \textit{scripts} auxiliares de ingestão, avaliação e registro dos resultados. Demonstrou-se, assim, que é possível estruturar uma base semântica consultável e responder perguntas em linguagem natural com apoio em evidências recuperadas, aproximando o uso de LLMs das exigências de atualidade, verificabilidade e controle das fontes apontadas pela literatura \parencite{gao2023ragSurvey}. A contribuição é, portanto, metodológica e aplicada, pois não se limita ao artefato implementado: envolve também um protocolo reprodutível, preservação de metadados legislativos e uma estratégia explícita de avaliação dos resultados.

Sob o enquadramento de \textit{Design Science Research} adotado no artigo, essa contribuição deve ser entendida como construção e avaliação inicial de um artefato de design para mediação informacional em acervo parlamentar, e não como implantação operacional validada em ambiente institucional.

\subsection{Discussão dos resultados frente aos objetivos}

À luz dos objetivos propostos, os resultados constituem evidência inicial de viabilidade. A prova de conceito atendeu aos critérios centrais definidos para indexação da base de discursos da 56\textordfeminine{} Legislatura, recuperação de trechos relevantes, geração de respostas fundamentadas e avaliação estruturada. As três rodadas automatizadas completaram a bateria de 20 perguntas sem falhas operacionais e apresentaram médias entre 9,15 e 9,55 em escala de 0 a 10. A inspeção manual e a triangulação amostral reforçaram, contudo, que a qualidade de um sistema RAG institucional depende de disciplina metodológica e de decisões arquiteturais consistentes. Estratégias de \textit{chunking}, desenho do \textit{prompt}, qualidade dos metadados, escolha do LLM juiz e avaliação iterativa influenciam diretamente o comportamento do \textit{chatbot}.

Tal constatação é compatível com a literatura, segundo a qual RAG não deve ser entendido apenas como a conexão entre um LLM e um banco vetorial, mas como uma solução sociotécnica que exige avaliação contínua, rastreabilidade, governança e critérios explícitos de sucesso, especialmente quando utilizada em funções públicas sensíveis \parencite{gao2023ragSurvey,deAlmeidaSantos2025aiGovernance,saadFalconEtAl2024ares}. Nesse sentido, a principal contribuição do trabalho não reside apenas na implementação do protótipo, mas também na explicitação das escolhas técnicas, limitações práticas e condições institucionais envolvidas em sua construção.

\subsection{Limitações}

Entre as limitações atuais da solução, destacam-se o tamanho ainda reduzido da bateria automatizada, a necessidade de validação humana mais ampla, a ausência, nesta etapa, de filtros mais sofisticados por autor, período ou tipo de discurso diretamente integrados ao fluxo de geração, a possibilidade de respostas excessivamente amplas em consultas abertas e a dependência de serviços externos para embeddings e parte da geração. Além disso, testes com LLM juiz alternativo indicaram que a estratégia de \textit{LLM as a Judge} é sensível ao modelo escolhido, podendo produzir avaliações excessivamente permissivas. Tais limitações dialogam com alertas da literatura acerca de riscos operacionais, vieses de dados, segmentação documental e necessidade de monitoramento contínuo em sistemas RAG \parencite{gao2023ragSurvey,wangEtAl2025segmentation,liuEtAl2025judgeRag}.

Há também limitações institucionais que não se resolvem apenas por ajuste técnico do pipeline. A qualidade de um sistema desse tipo depende de curadoria contínua do acervo, atualização controlada da base, definição de responsáveis por validação documental e regras para lidar com registros corrigidos, removidos ou incompletos. Por isso, os resultados devem ser lidos como evidência de viabilidade para consulta assistida e pesquisa aplicada, não como autorização para uso autônomo em processos decisórios ou atendimento público sem governança adicional.

\subsection{Trabalhos futuros}

Como trabalhos futuros, recomenda-se ampliar a bateria de avaliação, balancear melhor as categorias de perguntas, comparar avaliações por LLM como juiz com validações humanas independentes e experimentar \textit{frameworks} especializados, como RAGAS, para complementar a rubrica utilizada na avaliação por LLM. Recomenda-se, ainda, incorporar filtros estruturados por metadado na etapa de recuperação, experimentar reranqueadores mais robustos, testar diferentes estratégias de \textit{chunking} por tipo documental e implementar pós-processamento determinístico para a seção de referências.

Além disso, a prova de conceito pode ser adaptada para outros tipos de documentos legislativos, como pareceres, projetos de lei, notas técnicas e documentos administrativos relacionados ao processo legislativo. Essa expansão exigiria reavaliar metadados, estratégias de \textit{chunking}, categorias de pergunta, rubrica e critérios de validação em função de cada novo acervo. Do ponto de vista institucional, trabalhos futuros também devem detalhar protocolos de atualização da base, papéis de curadoria, validação humana, monitoramento de erros, registro de interações e critérios para separar respostas exploratórias de respostas aptas a circular em fluxos formais. Em horizonte mais avançado, um desdobramento possível consiste na evolução do sistema para cenários de geração assistida de minutas, relatórios ou discursos, desde que preservados mecanismos de verificação documental, supervisão humana obrigatória e controles institucionais adequados. Dessa forma, o presente trabalho abre caminho para agenda mais ampla de pesquisa e desenvolvimento sobre inteligência artificial generativa responsável no contexto do Parlamento brasileiro.
